\chapter{Appendix}
\label{chap:appendix}

\section{Synthetic Regime Ground-Truth Validation: Full Derivation}
\label{app:synthetic-validation}

This appendix gives the full construction and calibration of the synthetic regime-switching process summarized in Section~\ref{subsec:regime-synthetic-validation}, used to measure regime-recovery accuracy against a known ground truth --- something no validation against real market data can provide, since the true underlying regime process, if one exists, is never observed there.

A synthetic regime-switching process was constructed and calibrated to the real, already-fitted five-asset model: a six-state Markov-modulated geometric Brownian motion using the fitted transition matrix and each state's empirical annualized return and volatility directly as the generating parameters, so that the synthetic process matches, by construction, what the fitted model itself believes the data-generating process looks like.

Two independent discretization schemes were implemented to simulate price paths from this process, so that any recovery result could be checked for sensitivity to the simulation method itself rather than trusting a single generator. Level-space Euler--Maruyama and Milstein schemes were both implemented for the underlying SDE, and an empirical convergence-order comparison was attempted between them, holding the regime fixed to isolate discretization error from regime-switching noise; at the hourly step size and volatility levels relevant to this project, the theoretically expected $O(\Delta t)$ discretization bias between the two schemes ($\approx 0.007$ in variance units) was found analytically to sit roughly 28 times below the Monte Carlo noise floor achievable at a feasible sample size (200{,}000 paths, noise floor $\approx 0.19$), so no resolvable difference between the two schemes was expected or found --- a legitimate methodological conclusion (Euler--Maruyama's multiplicative discretization is already highly accurate for GBM at these parameters) rather than an inconclusive test. A second, non-SDE generator was implemented independently: a regime-switching trinomial tree with Cox--Ross--Rubinstein-style fixed jump sizes per state, whose three jump probabilities are solved numerically (\texttt{scipy.optimize}, Nelder--Mead) to match each state's target mean and variance exactly, subject to summing to one --- a genuine small calibration problem rather than a closed-form shortcut.

Three further numerical cross-checks validated the calibrated process's implied distributional properties against independent methods before it was used for anything else. A Crank--Nicolson finite-difference solution of the Kolmogorov forward equation for log-price, regime held fixed, was checked against the closed-form Gaussian solution available in that special case; an initial implementation returned a variance exactly double the closed-form value on every tested state (a coefficient error --- the Crank--Nicolson $\theta=0.5$ averaging factor had been omitted from the diffusion coefficient), and once corrected, the finite-difference solution matched the closed form to 0.04\% relative error uniformly across all six states. The regime-weighted mixture density's tail probabilities were computed by direct numerical integration (\texttt{scipy.integrate.quad}) as a semi-analytical alternative to simulation. Finally, the exact discrete-time characteristic function of the cumulative return over a fixed horizon was derived via a transfer-matrix recursion over the six states (not an approximation) and inverted by FFT to obtain the model-implied 24-hour return density; the resulting 24-hour variance (0.00103) came out roughly half of what treating each hour's regime draw as independent would predict (0.00206), consistent with the fitted transition matrix's high self-persistence ($\sim$93--96\% per hour) making a 24-hour window far more likely to stay within one regime than an i.i.d.\ approximation would suggest.

With the generator validated, the central question --- does the actual feature-engineering and HMM-fitting pipeline (Sections~\ref{subsec:pipeline-stages}, \ref{subsec:regime-model-selection}--\ref{subsec:regime-final-fit}) recover the true regime label from data it did not see the labels for --- was tested by simulating synthetic OHLCV series with a known regime at every hour, running them through the unmodified production feature pipeline and a fresh single-asset HMM fit, and scoring the recovered labels against the known truth with Adjusted Rand Index and a Hungarian-matched classification accuracy (\texttt{scipy.optimize.linear\_sum\_assignment} on the confusion matrix, which is invariant to the arbitrary label ordering EM produces). Five candidate explanations for the resulting recovery level were tested directly rather than assumed, each isolating one variable: more calendar time per path (one year versus three, single asset), the actual best-of-five-seed production fitting procedure in place of a single seed, pooling multiple statistically independent synthetic assets, pooling multiple assets with a strongly and directly verified cross-asset regime correlation (a shared latent market-regime process each asset's own regime follows with high, controlled probability), and matching production's real training-set scale exactly ($\approx$246{,}000 pooled rows). None of the five moved recovery accuracy outside a narrow band (Adjusted Rand Index 0.15--0.20 across every configuration tested, including the one built to match production's real scale and asset count directly), ruling out data quantity, cross-asset correlation, and seed count as explanations in this setting.

The eventual explanation was obtained directly rather than inferred: a Bayes-optimal classifier constructed from the \emph{true} generating parameters --- no feature engineering, no HMM estimation, the best possible single-observation classification given perfect knowledge of the answer --- achieved an Adjusted Rand Index of only 0.047, below what the full estimated pipeline achieved in every configuration tested. Inspecting the calibrated parameters directly explains why: the six states' hourly log-return means are all close to zero and close to one another (within an order of magnitude), and even the volatility differences that do separate the states overlap substantially when only one hourly observation is available. The states are differentiated almost entirely by volatility level, and distinguishing six heavily overlapping, near-zero-mean Gaussians from single noisy draws is close to an irreducible classification problem regardless of how well its parameters are estimated. This also resolves an apparent tension with the seed-sensitivity results of Section~\ref{subsec:regime-engine-audit}: seed agreement measures whether independent EM restarts converge to the \emph{same estimated parameters}, which more data does improve, while the synthetic recovery test measures whether the resulting model correctly labels each individual hour, which is bounded by class overlap and is a different property that more data does not fix. A model can be perfectly reproducible across refits while remaining noisy at the per-timestep level; the two findings are consistent rather than contradictory once the two questions are separated.

\section{Implementation Note: A Multiprocessing Deadlock}
\label{app:blas-deadlock}

While parallelizing the five-seed HMM refit of Section~\ref{subsec:regime-engine-audit} across CPU cores, a silent hang --- not an incorrect result --- was traced to Python's default \texttt{fork}-based process pool deadlocking completely when launched immediately after a BIC-search fit large enough to trigger multi-threaded BLAS execution in the parent process. Forking while another thread holds a BLAS-internal lock duplicates only the calling thread, so the lock is never released in the child and the first BLAS call inside any worker hangs indefinitely; five worker processes were observed sitting at under 1\% CPU utilization for over fifteen minutes before the defect was diagnosed. Switching the process pool to the \texttt{spawn} start method, which gives every worker a genuinely fresh interpreter with no inherited thread state, resolved it; the fix was verified by confirming all five workers reached the expected 97--98\% CPU utilization and by the synthetic sanity check of Appendix~\ref{app:synthetic-validation} run on pure random noise, which correctly returned near-chance seed agreement once the pool was actually computing.
